\section{Limitations and Conclusion}\label{sec:conclusion}

\subsection{Threats to Validity}\label{sec:threats}
\emph{Internal validity.} The survey is cross-sectional and self-reported, so temporal order cannot be established, reverse causation (procrastinators turning to GenAI near deadlines) is plausible, and both measures share one questionnaire, which may inflate their correlation~\cite{podsakoff2003common}. Only age, gender and education were available as covariates; workload, stress and self-regulation remain unmeasured confounders. \emph{Statistical conclusion validity.} With $n=285$, the test had about 80\% power to detect $\rho_s\approx0.17$, and the three group contrasts were Holm-corrected; the sensitivity samples are not independent replications, because they share most respondents. \emph{Construct validity.} Exposure is a single frequency item that ignores the purpose of use, and the UPS measures unintentional rather than academic procrastination~\cite{fernie2017ups}. \emph{External validity.} The data are a convenience sample of Indian students collected for a study of academic dishonesty~\cite{sharma2026dataset}; the estimate does not transfer directly to Kazakhstan. The latency figures come from one two-core machine and from 1{,}000 rather than 10{,}000 resamples per call; absolute times will differ elsewhere, although the guardrail and agreement results should not depend on hardware.

\subsection{Conclusion and Future Work}
Subject to the threats in Section~\ref{sec:threats}, the answers are as follows. For RQ1, academic GenAI-use frequency was positively and robustly associated with unintentional procrastination ($\hat\rho_s=\vrhoMain$), and the estimate was stable under two alternative sample selections (\vrhoA{} and \vrhoB). For RQ2, the association is modest: frequent users scored 3.5--4 points higher than the lowest-use group, and frequency added about 8\% of explained variance to a demographic model. For RQ3, vectorising the bootstrap reduced median latency by a factor of \vspeedMin--\vspeedMax{} while reproducing every estimate.

Future work should (i)~validate Kazakh and Russian versions of the UPS, (ii)~follow students in Kazakhstan over three waves to separate the temporal order of use and procrastination, and (iii)~measure the purpose of GenAI use and behavioural deadline data alongside frequency. The code, data derivatives and validation record are available at \url{https://github.com/yernurkussain-source/TASK2_KUSSAIN_ZHAHYP}.
